\begin{center}
    {\Large\itshape Alban au XV\textsuperscript{ème} siècle (suite)}
\end{center}

\bigskip

En remontant dans le temps, nous aurions pu assister au mariage de noble Raymond de Peyrole d'Ambialet, avec Giscarde de la Roque, célébré le 30 janvier 1461 dans l'église du Travet (Bertrand de la Roque étant le seigneur de ce fief)... L'acte de mariage ne fut enregistré par le notaire de Valence, Maître Bénet que le 7 février 1475. Ce qui nous y intéresse, c'est qu'un albanais, Dominique de Moret, qui s'est enrichi dans le commerce du fer (rappelez-vous qu'il y a une mouline ferran), est témoin à la signature du contrat et que le notaire écrit à son sujet en latin D. de Moréto de Albanio (le mot est orthographié Alba NHO). Et ceci pourrait bien sembler un écho de l'accord passé deux siècles avant entre l'évêque d'Albi et le Père Abbé de Saint Benoît de Castres où il est effectivement question de Ecclesia Sanctae Marie de ALBANIO, et confirmerait d'une certaine façon l'hypothèse émise de l'appellation non pas de Notre Dame d'Alban mais de Notre Dame « des ALBANS » (les Albans étant une tribu gauloise qui aurait essaimé quelques siècles avant Jésus Christ et aurait fondé sur les rivages de l'Adriatique l'actuelle République d'Albanie). Ceci écrit pour le panache des actuels Albanais d'Alban. À cette date de 1475, les de Moret sont coseigneurs de Nougaret avec Bertrand de St Pol de Bonneval.
